
\documentclass[UTF8,a4paper,12pt]{ctexart}
\usepackage{geometry}
\geometry{left=2.45cm,right=2.45cm,top=2.55cm,bottom=2.45cm}
\usepackage{amsmath,amssymb,bm,graphicx,booktabs,longtable,float,array,caption,fancyhdr,setspace,tocloft,enumitem,hyperref}
\setstretch{1.45}
\setlength{\headheight}{15pt}
\pagestyle{fancy}\fancyhf{}\lhead{2017A CT系统参数标定及成像应用}\rhead{返工版}\cfoot{\thepage}
\captionsetup{font=small,labelsep=quad}
\setcounter{tocdepth}{1}
\renewcommand{\contentsname}{目录}
\begin{document}
\begin{center}
{\heiti\zihao{2} 基于逐投影几何校正与滤波反投影的CT系统参数标定及成像应用}\par
\vspace{0.6cm}
{\zihao{4} 2017A 数学建模返工版}\par
\end{center}

\begin{abstract}
针对2017年高教社杯全国大学生数学建模竞赛A题，本文重新建立了可验证的CT系统参数标定与成像模型。旧的直接反投影方法没有让标定参数进入重建链，模板回代误差较大；返工版首先以附件1模板为已知吸收率函数，计算其Radon理论投影，再以附件2实测投影为目标，联合标定角度零点、角度步长和探测器尺度，并对每一列投影估计独立的探测器中心偏移，形成“角度--尺度--逐投影shift”的校正模型。标定结果为：角度序列
\(\theta_j=30.0897^\circ+(j-1)\times 0.9958^\circ\)，探测器尺度为1.4147 cell/radon-pixel，等价探测器间距为0.7069像素，即0.2761 mm；模板投影匹配平均相关系数达到0.9987。

在该标定参数下，本文将附件2、3、5的512探测器投影统一映射回模板Radon坐标，再采用Shepp--Logan滤波反投影重建。模板回代验证显示，baseline直接FBP的RMSE为0.3735、SSIM为0.0853；返工主模型的MAE为0.0279、RMSE为0.0607、SSIM为0.6826，投影回代RMSE为1.0017、相关系数为0.9995，说明标定参数已真实进入重建并显著改善结果。问题二10个指定点吸收率依次为0.0133, 0.8420, 0.1695, 1.5202, 1.0584, 1.6634, 1.0873, 0.0302, 0.0000, 0.0194；问题三10个指定点吸收率依次为1.1142, 2.0990, 5.5972, 0.6480, 0.9869, 2.5076, 5.3018, 0.8477, 7.2589, 0.0922。本文同时输出返工版\texttt{problem2\_rework.xls}、\texttt{problem3\_rework.xls}和支撑材料中的完整xlsx矩阵。

针对标定稳定性，本文对初始角和探测器尺度进行扰动，报告RMSE和SSIM变化，并设计由非中心圆盘、环形薄边、多方向条纹和灰度阶梯块组成的新模板。模拟投影剖面距离显示，新模板可提高不同角度投影的可区分度，从而改善角度和尺度标定的可辨识性。
\end{abstract}
\noindent\textbf{关键词：} CT参数标定；Radon变换；逐投影shift校正；滤波反投影；投影回代验证
\clearpage
\tableofcontents
\clearpage

\section{问题重述}
CT系统通过射线穿过介质后的衰减信息恢复内部结构。二维平行束CT中，一个角度下探测器阵列记录的是吸收率函数沿一组平行直线的积分，多个角度的投影共同构成sinogram。若旋转中心、探测器间距和射线方向准确，则可以通过反投影或迭代重建获得介质吸收率；若安装误差存在，则图像会出现位移、尺度变形和伪影。

本题给出一个已知模板的吸收率矩阵附件1及其接收信息附件2，要求标定CT系统旋转中心、探测器单元间距和180个射线方向。随后利用同一系统得到的附件3、附件5接收信息重建两个未知介质，并给出附件4中10个位置处的吸收率。最后需要分析标定精度和稳定性，并自行设计新模板以改进标定质量。

\section{问题分析}
本题的关键不是单纯使用库函数做反投影，而是让由模板得到的系统几何参数真实进入未知介质重建。若直接把512×180接收矩阵送入默认\texttt{iradon}，等价于假设探测器中心、尺度和角度均理想，这会忽略题目要求标定的核心参数。返工版因此将问题拆为三层：第一层，用模板Radon理论投影解释附件2；第二层，利用标定得到的角度、尺度和逐投影shift校正附件3、附件5；第三层，通过模板回代和未知投影回代检验结果。

问题一属于参数估计问题。本文将待估参数分为全局参数和逐投影参数：全局参数包括角度零点\(\theta_0\)、角度步长\(\Delta\theta\)、探测器尺度\(\lambda\)；逐投影参数为每列接收信息的探测器中心偏移\(b_j\)。问题二、三属于图像反演问题，在标定后的统一Radon坐标中用滤波反投影求解。问题四属于稳定性与实验设计问题，需要同时考察模型误差、参数扰动和模板可辨识性。

\section{数据审计、假设与符号}
附件1为256×256模板吸收率矩阵；附件2、附件3、附件5均为512×180接收矩阵；附件4为10个物理坐标点。托盘边长按题图单位取100 mm，因此像素尺寸为\(100/256=0.390625\) mm。本文保留原始数据，所有返工输出写入\texttt{支撑材料\_返工版}，并将关键数字冻结到\texttt{frozen\_numbers\_rework.json}。

基本假设如下：1）系统为平行束CT，探测器等距排列；2）模板和未知介质在扫描过程中固定不动；3）接收值与线积分经线性增益和偏置近似相关；4）附件2、3、5由同一套系统几何参数产生，因此未知介质不得重新调参；5）FBP产生的少量负值为数值振铃，最终吸收率取非负。

\begin{table}[H]\centering\caption{主要符号说明}\begin{tabular}{lll}\toprule
符号&含义&说明\\\midrule
\(f(x,y)\)&吸收率分布&二维图像函数\\
\(p_j(s)\)&第\(j\)个角度的投影&512探测器接收序列\\
\(\theta_0,\Delta\theta\)&角度零点与步长&全局标定参数\\
\(\lambda\)&探测器尺度&cell/radon-pixel\\
\(b_j\)&第\(j\)列探测器中心偏移&逐投影校正参数\\
\(\hat f\)&重建吸收率矩阵&256×256\\\bottomrule
\end{tabular}\end{table}

\section{问题一模型建立与参数标定}
\subsection{Radon投影模型}
设旋转中心坐标系下的射线方程为
\begin{equation}
(x-x_0)\cos\theta_j+(y-y_0)\sin\theta_j=s_i,
\end{equation}
则理想投影为
\begin{equation}
p_j(s)=R_{\theta_j}f(s)=\iint f(x,y)\delta((x-x_0)\cos\theta_j+(y-y_0)\sin\theta_j-s)\,dxdy.
\end{equation}
附件1模板\(f_0\)已知，因此可以计算任意候选角度下的理论Radon剖面。由于题目探测器有512个单元，而256×256模板Radon变换的自然detector长度为363，必须标定二者坐标映射。

\subsection{角度--尺度--逐投影shift联合标定}
令第\(j\)列角度为
\begin{equation}
\theta_j=\theta_0+j\Delta\theta,\quad j=0,1,\ldots,179.
\end{equation}
若模板Radon坐标为\(u\)，512探测器中心化坐标为\(c_i=i-255.5\)，则映射关系写成
\begin{equation}
c_i=\lambda u+b_j.
\end{equation}
给定\(\theta_0,\Delta\theta,\lambda\)后，对每一列用互相关求最优\(b_j\)，目标为最大化平均标准化相关系数：
\begin{equation}
\max_{\theta_0,\Delta\theta,\lambda,\{b_j\}}\frac1{180}\sum_{j=0}^{179}\rho\left(y_j,\,\mathcal I_{\lambda,b_j}(R_{\theta_j}f_0)\right).
\end{equation}
本文先做粗网格搜索，再用Nelder--Mead局部优化，最后对180列逐列求shift。得到\(\theta_0=30.0897^\circ\)、\(\Delta\theta=0.9958^\circ\)、\(\lambda=1.4147\)，平均相关系数为0.9987。探测器间距为\(1/\lambda=0.7069\)像素，即0.2761 mm。

\begin{figure}[H]\centering\includegraphics[width=.82\textwidth]{figs/angle_sequence_rework.png}\caption{返工版180个射线方向标定结果}\end{figure}
\begin{figure}[H]\centering\includegraphics[width=.82\textwidth]{figs/detector_shift_sequence.png}\caption{逐投影探测器中心偏移校正序列}\end{figure}
图1说明角度序列近似线性递增但零点约为30度；图2说明探测器偏移并非常数，其均值为-22.5625 cell、标准差为17.3106 cell。返工版不再把中心偏移压缩为一个未使用的旋转中心数字，而是直接把\(b_j\)用于每列sinogram校正。

\subsection{标定结果局部明细}
\begin{table}[H]\centering\caption{前24列投影匹配结果}\resizebox{\textwidth}{!}{\begin{tabular}{ccccc}
\toprule
projection\_index & theta\_deg & scale\_cell\_per\_radon\_pixel & shift\_cell & corr \\
\midrule
0.0000 & 30.0897 & 1.4147 & 18.2500 & 0.9999 \\
1.0000 & 31.0855 & 1.4147 & 17.2500 & 1.0000 \\
2.0000 & 32.0813 & 1.4147 & 17.0000 & 0.9998 \\
3.0000 & 33.0770 & 1.4147 & 16.2500 & 0.9999 \\
4.0000 & 34.0728 & 1.4147 & 15.5000 & 0.9999 \\
5.0000 & 35.0686 & 1.4147 & 15.0000 & 0.9999 \\
6.0000 & 36.0644 & 1.4147 & 14.2500 & 0.9999 \\
7.0000 & 37.0601 & 1.4147 & 13.5000 & 0.9999 \\
8.0000 & 38.0559 & 1.4147 & 13.0000 & 0.9999 \\
9.0000 & 39.0517 & 1.4147 & 12.0000 & 0.9999 \\
10.0000 & 40.0475 & 1.4147 & 11.5000 & 0.9999 \\
11.0000 & 41.0433 & 1.4147 & 10.7500 & 0.9999 \\
12.0000 & 42.0390 & 1.4147 & 10.0000 & 0.9999 \\
13.0000 & 43.0348 & 1.4147 & 9.0000 & 0.9998 \\
14.0000 & 44.0306 & 1.4147 & 7.0000 & 0.9991 \\
15.0000 & 45.0264 & 1.4147 & 7.0000 & 0.9997 \\
16.0000 & 46.0221 & 1.4147 & 6.0000 & 0.9994 \\
17.0000 & 47.0179 & 1.4147 & 6.0000 & 0.9998 \\
18.0000 & 48.0137 & 1.4147 & 5.0000 & 0.9997 \\
19.0000 & 49.0095 & 1.4147 & 4.0000 & 0.9995 \\
20.0000 & 50.0053 & 1.4147 & 4.0000 & 0.9999 \\
21.0000 & 51.0010 & 1.4147 & 3.7500 & 1.0000 \\
22.0000 & 51.9968 & 1.4147 & 3.0000 & 0.9999 \\
23.0000 & 52.9926 & 1.4147 & 2.0000 & 0.9999 \\
\bottomrule
\end{tabular}}\end{table}
表2显示前若干列匹配相关系数接近1，说明模板理论投影能够较好解释附件2。与旧版不同，返工版保存完整180列匹配表，后续未知介质重建直接复用这些参数。

\section{问题二和问题三的重建模型}
\subsection{校正sinogram}
对任意原始接收矩阵\(Y\)，返工版按标定参数将512探测器数据映射回模板Radon坐标：
\begin{equation}
\tilde Y(u,j)=Y(\lambda u+b_j,j).
\end{equation}
该式是本次返工的核心：角度、尺度和逐投影shift都进入实际数据处理链。之后对\(\tilde Y\)使用Shepp--Logan滤波反投影：
\begin{equation}
\hat f(x,y)=\int_0^\pi [\tilde p_\theta(s)*h_{SL}(s)]_{s=x\cos\theta+y\sin\theta}\,d\theta.
\end{equation}
幅值标定使用模板回代重建与附件1之间的线性回归，得到统一的\(a,b\)，未知介质统一使用同一幅值映射并取非负。

\subsection{问题二结果}
\begin{figure}[H]\centering\includegraphics[width=.72\textwidth]{figs/problem2_reconstruction_rework.png}\caption{返工版问题二未知介质重建图}\end{figure}
问题二重建图显示介质主要位于中上部及中部区域，低吸收背景被明显压低。与旧版相比，返工版减少了“整幅图大面积非零”的扩散现象。主要连通组件见表3。
\begin{table}[H]\centering\caption{问题二主要连通组件几何信息}\resizebox{\textwidth}{!}{\begin{tabular}{ccccccccc}
\toprule
area\_mm2 & centroid\_x\_mm & centroid\_y\_mm & mean\_absorption & max\_absorption & bbox\_xmin\_mm & bbox\_xmax\_mm & bbox\_ymin\_mm & bbox\_ymax\_mm \\
\midrule
2407.5317 & 49.6635 & 54.3349 & 1.0059 & 1.8448 & 27.9297 & 71.6797 & 10.7422 & 91.9922 \\
\bottomrule
\end{tabular}}\end{table}

\subsection{问题三结果}
\begin{figure}[H]\centering\includegraphics[width=.72\textwidth]{figs/problem3_reconstruction_rework.png}\caption{返工版问题三未知介质重建图}\end{figure}
问题三投影强度明显高于问题二，重建结果呈现多个高吸收区域。由于问题三吸收率范围较宽，本文在解释时同时报告点值、连通域均值和投影回代残差，而不只依据视觉图像判断形状。
\begin{table}[H]\centering\caption{问题三主要连通组件几何信息}\resizebox{\textwidth}{!}{\begin{tabular}{ccccccccc}
\toprule
area\_mm2 & centroid\_x\_mm & centroid\_y\_mm & mean\_absorption & max\_absorption & bbox\_xmin\_mm & bbox\_xmax\_mm & bbox\_ymin\_mm & bbox\_ymax\_mm \\
\midrule
2002.8687 & 49.0205 & 48.3518 & 2.9145 & 6.7524 & 4.4922 & 95.5078 & 8.3984 & 84.9609 \\
391.0828 & 76.1942 & 22.6466 & 3.0282 & 7.6872 & 62.3047 & 93.1641 & 5.2734 & 41.6016 \\
170.8984 & 20.7401 & 83.4992 & 3.2128 & 7.7000 & 4.8828 & 34.5703 & 73.6328 & 90.8203 \\
46.9971 & 13.6858 & 65.1697 & 2.9330 & 5.7309 & 8.7891 & 20.1172 & 61.9141 & 70.1172 \\
38.1470 & 79.6719 & 83.3469 & 2.5065 & 3.9456 & 74.0234 & 85.3516 & 78.3203 & 87.6953 \\
\bottomrule
\end{tabular}}\end{table}

\subsection{10个指定点吸收率}
\begin{table}[H]\centering\caption{附件4指定位置吸收率返工版结果}\begin{tabular}{ccccc}
\toprule
序号 & x/mm & y/mm & 问题2吸收率 & 问题3吸收率 \\
\midrule
1.0000 & 10.0000 & 18.0000 & 0.0133 & 1.1142 \\
2.0000 & 34.5000 & 25.0000 & 0.8420 & 2.0990 \\
3.0000 & 43.5000 & 33.0000 & 0.1695 & 5.5972 \\
4.0000 & 45.0000 & 75.5000 & 1.5202 & 0.6480 \\
5.0000 & 48.5000 & 55.5000 & 1.0584 & 0.9869 \\
6.0000 & 50.0000 & 75.5000 & 1.6634 & 2.5076 \\
7.0000 & 56.0000 & 76.5000 & 1.0873 & 5.3018 \\
8.0000 & 65.5000 & 37.0000 & 0.0302 & 0.8477 \\
9.0000 & 79.5000 & 18.0000 & 0.0000 & 7.2589 \\
10.0000 & 98.5000 & 43.5000 & 0.0194 & 0.0922 \\
\bottomrule
\end{tabular}\end{table}
表5为题目第二、三问指定点的直接数值答案。问题二中第4、5、6、7点较高，说明这些点位于主体介质内；第1、9点接近背景。问题三中第3、7、9点较高，体现其内部吸收率非均匀性更强。

\section{模型检验与对比}
\subsection{模板重建对比}
\begin{table}[H]\centering\caption{模板重建误差：baseline与返工主模型对比}\resizebox{\textwidth}{!}{\begin{tabular}{cccccccc}
\toprule
name & amp\_a & amp\_b & mae & rmse & ssim & min & max \\
\midrule
main\_corrected\_shepp-logan & 1.4348 & -0.0000 & 0.0279 & 0.0607 & 0.6826 & -0.5000 & 1.0106 \\
main\_corrected\_ramp & 1.4326 & 0.0003 & 0.0291 & 0.0610 & 0.6461 & -0.5165 & 1.0202 \\
main\_corrected\_hann & 1.4408 & -0.0008 & 0.0264 & 0.0623 & 0.7447 & -0.4219 & 1.0068 \\
baseline\_raw512\_hann & 0.5417 & 0.1011 & 0.2788 & 0.3734 & 0.0862 & 0.0766 & 0.3700 \\
baseline\_raw512\_shepp-logan & 0.5375 & 0.1018 & 0.2791 & 0.3735 & 0.0853 & 0.0658 & 0.3723 \\
baseline\_raw512\_ramp & 0.5361 & 0.1020 & 0.2791 & 0.3736 & 0.0848 & 0.0546 & 0.3790 \\
\bottomrule
\end{tabular}}\end{table}
表6表明，直接512探测器FBP的RMSE约为0.3735、SSIM约为0.0853；返工主模型的RMSE降至0.0607、SSIM提升至0.6826。这说明几何校正不是论文装饰，而是真正改善了模板重建。
\begin{figure}[H]\centering\includegraphics[width=.78\textwidth]{figs/baseline_comparison_rmse.png}\caption{baseline与返工模型模板RMSE对比}\end{figure}
\begin{figure}[H]\centering\includegraphics[width=.72\textwidth]{figs/template_reconstruction_rework.png}\caption{返工版模板回代重建图}\end{figure}

\subsection{投影回代检验}
\begin{table}[H]\centering\caption{附件2模板投影回代误差}\resizebox{\textwidth}{!}{\begin{tabular}{cccc}
\toprule
name & proj\_rmse & proj\_mae & proj\_corr \\
\midrule
main\_corrected\_ramp & 0.9952 & 0.5745 & 0.9995 \\
main\_corrected\_shepp-logan & 1.0017 & 0.5746 & 0.9995 \\
main\_corrected\_hann & 1.0382 & 0.5928 & 0.9995 \\
baseline\_raw512\_hann & 20.2614 & 13.9544 & 0.7892 \\
baseline\_raw512\_shepp-logan & 20.2807 & 13.9804 & 0.7887 \\
baseline\_raw512\_ramp & 20.2871 & 13.9893 & 0.7886 \\
\bottomrule
\end{tabular}}\end{table}
主模型投影回代RMSE为1.0017，MAE为0.5746，相关系数为0.9995。投影回代检验比单纯图像相似度更直接，因为它检查\(A\hat f\)是否能重新解释原始接收信息。

\begin{table}[H]\centering\caption{未知介质投影回代误差}\begin{tabular}{cccccc}
\toprule
name & proj\_rmse & proj\_mae & proj\_corr & proj\_scale & proj\_offset \\
\midrule
problem2\_rework & 4.4702 & 2.6914 & 0.9939 & 0.6898 & -1.9120 \\
problem3\_rework & 8.1692 & 5.8811 & 0.9975 & 0.6803 & -4.6802 \\
\bottomrule
\end{tabular}\end{table}
表8说明，问题二、三重建图重新投影后与原始接收矩阵具有较高相关系数。问题三RMSE较大，主要是因为其投影值尺度更高、内部高吸收区域更强，后续可用TV正则或显式系统矩阵迭代法改进。

\subsection{灵敏度分析}
\begin{table}[H]\centering\caption{返工版参数扰动灵敏度分析}\begin{tabular}{cccc}
\toprule
参数 & 扰动 & RMSE & SSIM \\
\midrule
初始角/度 & -0.5000 & 0.0696 & 0.6612 \\
初始角/度 & -0.2500 & 0.0640 & 0.6800 \\
初始角/度 & 0.0000 & 0.0607 & 0.6826 \\
初始角/度 & 0.2500 & 0.0602 & 0.6854 \\
初始角/度 & 0.5000 & 0.0626 & 0.6796 \\
探测器尺度相对扰动 & -0.0400 & 0.1280 & 0.6724 \\
探测器尺度相对扰动 & -0.0200 & 0.0930 & 0.6906 \\
探测器尺度相对扰动 & 0.0000 & 0.0607 & 0.6826 \\
探测器尺度相对扰动 & 0.0200 & 0.0832 & 0.5870 \\
探测器尺度相对扰动 & 0.0400 & 0.1190 & 0.4494 \\
\bottomrule
\end{tabular}\end{table}
\begin{figure}[H]\centering\includegraphics[width=.78\textwidth]{figs/sensitivity_rework.png}\caption{返工版参数扰动RMSE曲线}\end{figure}
灵敏度结果用于判断参数附近的误差变化。与旧版不同，本文不把所有扰动都解释为“更稳定”；若扰动后误差接近或略有改善，说明模型仍存在局部非凸和离散采样误差，应在论文中诚实说明。这也是问题四需要改进模板的原因。

\section{问题四：新模板设计}
当前模板由少数均匀介质组成，部分角度投影相似，导致角度与shift估计仍有局部歧义。本文设计新模板：非中心大圆盘提供稳定强峰，环形薄边提供尖锐边缘，多方向矩形条降低角度歧义，灰度阶梯块分离几何误差与增益误差。
\begin{figure}[H]\centering\includegraphics[width=.70\textwidth]{figs/new_template_design.png}\caption{返工版新模板设计示意图}\end{figure}
为量化可辨识性，本文比较不同角度标准化投影剖面的两两距离。距离越大，说明不同角度越不容易混淆。
\begin{table}[H]\centering\caption{原模板与新模板投影可辨识性对比}\begin{tabular}{cccc}
\toprule
模板 & 最小角度剖面距离 & 5\%分位距离 & 平均距离 \\
\midrule
原模板 & 0.1026 & 2.2826 & 8.6776 \\
新设计模板 & 0.7292 & 5.0137 & 16.0723 \\
\bottomrule
\end{tabular}\end{table}
表10表明，新模板在低分位和平均投影剖面距离上具有更强区分度，可用于提高角度零点和探测器尺度估计的稳定性。

\section{模型评价与改进}
返工版的优点是：第一，标定参数真实进入sinogram校正与重建链；第二，给出了baseline对比、模板回代和投影回代三类证据；第三，所有关键结果均保存为CSV/JSON/xlsx/xls文件，便于复现。局限是：逐投影shift是经验校正序列，尚未完全等价于机械旋转中心的解析参数；FBP仍可能有振铃和边缘过冲；问题三高吸收区域的定量精度仍受幅值线性标定影响。

后续改进方向包括：构造显式系统矩阵，将\(\theta_0,\Delta\theta,\lambda,b_j\)纳入ART/SART或TV正则重建；用新模板重新采集或模拟投影，联合估计中心、尺度、角度和增益；对指定点吸收率使用bootstrap或噪声扰动给出置信区间。

\section{结论}
本文完成2017A返工版建模。主要结果为：1）射线方向\(\theta_j=30.0897^\circ+(j-1)0.9958^\circ\)，探测器间距0.7069像素即0.2761 mm，模板匹配平均相关系数0.9987；2）返工主模型模板RMSE由baseline的0.3735降至0.0607，SSIM由0.0853升至0.6826，投影回代相关系数0.9995；3）问题二和问题三的10点吸收率分别为0.0133, 0.8420, 0.1695, 1.5202, 1.0584, 1.6634, 1.0873, 0.0302, 0.0000, 0.0194和1.1142, 2.0990, 5.5972, 0.6480, 0.9869, 2.5076, 5.3018, 0.8477, 7.2589, 0.0922，完整矩阵已保存；4）新模板模拟显示多特征非对称结构可提高角度投影可辨识性。

\clearpage
\section*{参考文献}\addcontentsline{toc}{section}{参考文献}
\begin{enumerate}[label={[\arabic*]}]
\item Kak A C, Slaney M. Principles of Computerized Tomographic Imaging[M]. IEEE Press, 1988.
\item Natterer F. The Mathematics of Computerized Tomography[M]. SIAM, 2001.
\item scikit-image developers. Radon transform and inverse Radon transform documentation[EB/OL].
\item Gonzalez R C, Woods R E. Digital Image Processing[M]. Pearson, 2018.
\item 全国大学生数学建模竞赛组委会. 2017年高教社杯全国大学生数学建模竞赛A题[Z]. 2017.
\end{enumerate}

\appendix
\section{附录A：支撑材料说明}
返工版支撑材料位于\texttt{支撑材料\_返工版}。其中\texttt{code/main\_rework.py}为完整求解脚本，\texttt{paper/论文.tex}和\texttt{paper/论文.pdf}为正式论文源文件与PDF，\texttt{results/frozen\_numbers\_rework.json}为冻结数字，\texttt{results/problem2.xls}和\texttt{results/problem3.xls}为真实Excel格式结果矩阵，\texttt{verification/verify\_outputs.py}用于核验输出文件。

\section{附录B：算法流程}
1. 读取附件1--5并审计维度；2. 计算模板Radon理论投影；3. 粗网格搜索\(\theta_0,\Delta\theta,\lambda\)，每列用互相关求\(b_j\)；4. 将512探测器数据映射到363长度Radon坐标；5. 用Shepp--Logan FBP重建模板、问题二和问题三；6. 用模板回归幅值标定；7. 计算10点吸收率、连通域、投影回代误差和灵敏度；8. 生成论文、图表和压缩包。

\end{document}
